\documentclass[a4paper,landscape]{article}
\usepackage[utf8]{inputenc}
\usepackage[italian]{babel}
\usepackage{tikz}
\usetikzlibrary{positioning}
\usepackage{geometry}
\geometry{margin=0.3cm, paperwidth=42cm, paperheight=29.7cm}

\tikzset{
  class/.style={
    rectangle, draw=black, rounded corners,
    minimum width=3.5cm,
    minimum height=1.2cm,
    font=\footnotesize,
    align=left,
    text width=3.5cm,
    inner sep=0.3cm
  },
  relation/.style={-latex, thick},
  every node/.style={anchor=center}
}

\begin{document}

% ==============================
% Diagramma NON RISTRUTTURATO
% ==============================
\section*{Diagramma non ristrutturato}

\begin{center}
\begin{tikzpicture}[node distance=2.5cm and 3.5cm]

% Livello superiore - Persona come entità principale
\node[class] (persona) at (0,0) {
  \textbf{Persona}\\[-0.3em]\rule{\linewidth}{0.2pt}\\
  + nome : string\\
  + cognome : string\\
  + codicefiscale : string\\
  + indirizzo : string\\
  + telefono : string\\
  + email : string
};

% Livello secondo - Specializzazioni di Persona
\node[class] (cliente) at (-4,-3) {
  \textbf{Cliente}\\[-0.3em]\rule{\linewidth}{0.2pt}\\
  + codcliente : int
};

\node[class] (dipendente) at (4,-3) {
  \textbf{Dipendente}\\[-0.3em]\rule{\linewidth}{0.2pt}\\
  + coddipendente : int
};

% Livello terzo - Entità collegate a Cliente
\node[class] (tessera) at (-4,-6) {
  \textbf{Tessera}\\[-0.3em]\rule{\linewidth}{0.2pt}\\
  + codtessera : int\\
  + numeropunti : numeric\\
  + dataemissione : date\\
  + datascadenza : date\\
  + stato : {ATTIVA, SOSPESA, SCADUTA}
};

% Livello centrale - Ordine come entità di collegamento
\node[class] (ordine) at (0,-6) {
  \textbf{Ordine}\\[-0.3em]\rule{\linewidth}{0.2pt}\\
  + codordine : int\\
  + prezzototale : real\\
  + dataacquisto : date
};

% Livello quarto - ArticoliOrdine
\node[class] (articoli) at (0,-9) {
  \textbf{ArticoliOrdine}\\[-0.3em]\rule{\linewidth}{0.2pt}\\
  + prezzo : numeric\\
  + numeroarticoli : int
};

% Livello destro - Prodotto e sue specializzazioni
\node[class] (prodotto) at (8,-6) {
  \textbf{Prodotto}\\[-0.3em]\rule{\linewidth}{0.2pt}\\
  + codprodotto : int\\
  + nome : string\\
  + descrizione : string\\
  + prezzo : numeric\\
  + luogoprovenienza : string\\
  + scorta : int
};

% Specializzazioni di Prodotto - disposte in griglia
\node[class] (frutta) at (5,-9) {
  \textbf{Frutta}\\[-0.3em]\rule{\linewidth}{0.2pt}\\
  + dataraccolta : date
};

\node[class] (verdura) at (8,-9) {
  \textbf{Verdura}\\[-0.3em]\rule{\linewidth}{0.2pt}\\
  + dataraccolta : date
};

\node[class] (latticini) at (11,-9) {
  \textbf{Latticini}\\[-0.3em]\rule{\linewidth}{0.2pt}\\
  + datamungitura : date\\
  + dataproduzione : date\\
  + datascadenza : date
};

\node[class] (farinacei) at (5,-12) {
  \textbf{Farinacei}\\[-0.3em]\rule{\linewidth}{0.2pt}\\
  + glutine : boolean
};

\node[class] (uova) at (8,-12) {
  \textbf{Uova}\\[-0.3em]\rule{\linewidth}{0.2pt}\\
  + datascadenza : date
};

\node[class] (confezionati) at (11,-12) {
  \textbf{Confezionati}\\[-0.3em]\rule{\linewidth}{0.2pt}\\
  + datascadenza : date
};

% Relazioni principali - Ereditarietà
\draw[relation] (persona) -- (cliente) node[midway, left] {ISA};
\draw[relation] (persona) -- (dipendente) node[midway, right] {ISA};

% Relazioni Cliente
\draw[relation] (cliente) -- (tessera) node[midway, left] {possiede};
\draw[relation] (cliente) -- (ordine) node[midway, below left] {effettua};

% Relazioni Dipendente
\draw[relation] (dipendente) -- (ordine) node[midway, below right] {gestisce};

% Relazioni Ordine-ArticoliOrdine-Prodotto
\draw[relation] (ordine) -- (articoli) node[midway, left] {contiene};
\draw[relation] (articoli) -- (prodotto) node[midway, below] {riferisce};

% Generalizzazioni Prodotto - disposte a raggiera
\draw[relation] (prodotto) -- (frutta) node[midway, above left] {ISA};
\draw[relation] (prodotto) -- (verdura) node[midway, left] {ISA};
\draw[relation] (prodotto) -- (latticini) node[midway, above right] {ISA};
\draw[relation] (prodotto) -- (farinacei) node[midway, below left] {ISA};
\draw[relation] (prodotto) -- (uova) node[midway, below] {ISA};
\draw[relation] (prodotto) -- (confezionati) node[midway, below right] {ISA};

\end{tikzpicture}
\end{center}


% ==============================
% Diagramma RISTRUTTURATO
% ==============================
\section*{Diagramma ristrutturato}

\begin{center}
\begin{tikzpicture}[node distance=2.5cm and 4cm]

% Livello superiore - Cliente e Dipendente affiancati
\node[class] (cliente2) at (-3,0) {
  \textbf{Cliente}\\[-0.3em]\rule{\linewidth}{0.2pt}\\
  + codcliente : int\\
  + nome : string\\
  + cognome : string\\
  + codicefiscale : string\\
  + indirizzo : string\\
  + telefono : string\\
  + email : string
};

\node[class] (dipendente2) at (3,0) {
  \textbf{Dipendente}\\[-0.3em]\rule{\linewidth}{0.2pt}\\
  + coddipendente : int\\
  + nome : string\\
  + cognome : string\\
  + codicefiscale : string\\
  + indirizzo : string\\
  + telefono : string\\
  + email : string
};

% Livello secondo - Tessera e Ordine
\node[class] (tessera2) at (-6,-3.5) {
  \textbf{Tessera}\\[-0.3em]\rule{\linewidth}{0.2pt}\\
  + codtessera : int\\
  + numeropunti : numeric\\
  + dataemissione : date\\
  + datascadenza : date\\
  + stato : STATOTESSERA
};

\node[class] (ordine2) at (0,-3.5) {
  \textbf{Ordine}\\[-0.3em]\rule{\linewidth}{0.2pt}\\
  + codordine : int\\
  + prezzototale : real\\
  + dataacquisto : date
};

% Livello terzo - Prodotto a destra
\node[class] (prodotto2) at (6,-3.5) {
  \textbf{Prodotto}\\[-0.3em]\rule{\linewidth}{0.2pt}\\
  + codprodotto : int\\
  + nome : string\\
  + descrizione : string\\
  + prezzo : numeric\\
  + luogoprovenienza : string\\
  + dataraccolta : date\\
  + datamungitura : date\\
  + glutine : boolean\\
  + datascadenza : date\\
  + dataproduzione : date\\
  + categoria : {FRUTTA, VERDURA, LATTICINI, FARINACEI, UOVA, CONFEZIONATI}\\
  + scorta : int
};

% Livello inferiore - ArticoliOrdine al centro
\node[class] (articoli2) at (0,-7) {
  \textbf{ArticoliOrdine}\\[-0.3em]\rule{\linewidth}{0.2pt}\\
  + prezzo : numeric\\
  + numeroarticoli : int
};

% Relazioni con etichette descrittive
\draw[relation] (cliente2) -- (tessera2) node[midway, above] {possiede (1:1)};
\draw[relation] (cliente2) -- (ordine2) node[midway, above left] {effettua (1:N)};
\draw[relation] (dipendente2) -- (ordine2) node[midway, above right] {gestisce (1:N)};
\draw[relation] (ordine2) -- (articoli2) node[midway, left] {contiene (1:N)};
\draw[relation] (articoli2) -- (prodotto2) node[midway, above] {riferisce (N:1)};

\end{tikzpicture}
\end{center}

\end{document}
